% ---------------------------------------------------------------- 摘要（第一部分分册）
% 与 tex/abstract.tex 同源，仅供“第一部分”分册使用：
%   * 保留全报告定位一句与“第一部分的主要发现”四条；
%   * 删去“第二部分的主要发现”四条（其口径、样本与结论属于第二部分，见另册），
%     并去掉指向第二部分章节的两个 \ref（分册中该标签不存在，会印成 ??）。
\section*{摘要}
\addcontentsline{toc}{section}{摘要}

\begin{keybox}[报告要点]
\small
本报告以\textbf{申万行业分类（2021 版）}为统一框架，分为两大部分：
第一部分从供需关系出发分析农业各子行业的周期性，以\hl{猪周期为基准}
测算跨品种、跨行业的相关性与传导时滞；第二部分系统梳理
申万农林牧渔行业\textbf{全部 \num{104} 家 A 股上市公司}近三年的经营与融资画像。
\textbf{本册为第一部分}，对应第一章至第六章；第二部分（逐公司画像与融资需求）
为另一册。

\textbf{第一部分的主要发现}：
（1）本轮农业下行（2024---2025）源于\hl{供给恢复而非需求萎缩}——
玉米连续增产、谷物进口大幅收缩，生猪出栏在产能未出清时继续增长，
白羽肉鸡祖代更新量创历史新高；
（2）品种周期长度由\hl{产能调整的生物学时间}决定，从蛋鸡的 \num{6}---\num{12} 个月
到肉牛的 \num{24}---\num{44} 个月以上，谱系跨越 \num{7} 倍；
（3）在统一口径（月度收益率、同一时间窗口）下，
\hl{跨品种相关性的绝对值普遍低于 \num{0.24}}——花生 \num{+0.48}、粳米 \num{+0.29}、
豆一 \num{+0.26} 三个品种略高于 5\% 显著性阈值，玉米 \num{+0.09}、豆粕 \num{+0.13}
等饲料原料品种均不显著，“猪价带动一切农产品”并非统计事实；
（4）截至 2026 年 9 月，生猪自繁自养已连续六个月处于亏损区间
（猪粮比价月度均值 \num{3.96}---\num{4.47}）、能繁母猪存栏同比转负，
而肉牛价格、蛋鸡存栏与水产投苗已率先反转，\hl{农业周期正处于
“猪链磨底、禽链与牛链上行”的分化阶段}。
\end{keybox}

\vspace{6pt}
\noindent\textbf{关键词}\quad 农业周期；猪周期；跨品种相关性；申万行业分类
\clearpage
